\chapter{The honest verdict}

\section{Is the project worth the effort?}

Yes, if the goal is a serious privacy-review assistant and an end-to-end computer-vision research project. No, if the goal is to claim that a single model can automatically make any photograph safe to share.

The project is technically worthwhile because it combines several real problems: instance segmentation, small-object detection, OCR, handwriting localization, dataset engineering, calibration, human-computer interaction, privacy-preserving export, and assurance. It is also useful because ordinary content-moderation scores do not directly answer the privacy question: \emph{what sensitive pixels were left visible?}

\begin{truthbox}[One-sentence outcome]
ConsentGuard V1 should examine a still image locally, collect high-recall evidence from specialist detectors, show every uncertain region to a human, destructively redact the approved regions, remove metadata, attack-test the result, and block export whenever a required check is missing or uncertain.
\end{truthbox}

\section{Why the current score is both useful and insufficient}

The frozen comparator has COCO-style mask mAP of \textbf{0.2335}. That is a respectable research baseline for nine heavily imbalanced privacy classes and a 4 GB training environment. It proves that the data loader, model, evaluation code, checkpointing, and main segmentation idea work.

It does not prove privacy safety. At the original score threshold of 0.5, the model covers \textbf{76.39\%} of labelled sensitive pixels, leaving \textbf{23.61\%} uncovered. The aggregate is dominated by large, common classes such as people. Handwriting recall is only \textbf{2.68\%}, licence-plate recall \textbf{13.57\%}, medicine recall \textbf{15.39\%}, and signature recall \textbf{28.92\%} at the same operating point.

\begin{warningbox}[The main lesson]
A model can improve mAP while still being unsafe for privacy. Conversely, lowering thresholds can recover sensitive regions while making the image unusable through over-redaction. V1 therefore needs class-specific thresholds, specialist models, human review, and leakage metrics.
\end{warningbox}

\section{The product boundary}

\begin{tabularx}{\textwidth}{@{}YY@{}}
\toprule
\textbf{Inside V1} & \textbf{Outside V1} \\
\midrule
Still images; local-first analysis; high-recall proposals; human correction; versioned audit; destructive redaction; metadata cleaning; independent assurance. & Video; biometric recognition; identity matching; pixel-based consent inference; automatic no-review release; legal-compliance guarantees; a promise of zero privacy risk. \\
\bottomrule
\end{tabularx}

The project should be described as a \textbf{review-assisted privacy system}. The output of the current Mask R-CNN is evidence supplied to that system, not the final product.

\section{What ``final model'' should mean}

There should be no single file named ``the final model.'' The release unit is a bundle:

\begin{itemize}[leftmargin=6mm]
  \item provider names and exact versions;
  \item trained checkpoint hashes;
  \item class ontology and geometry conventions;
  \item frozen threshold profile;
  \item evidence-fusion version;
  \item release-policy version;
  \item assurance configuration; and
  \item validation and locked-test reports.
\end{itemize}

The bundle becomes \file{release-candidate-v1} only after every gate in Chapter 11 passes. The complete reviewed system may then be tagged \file{ConsentGuard v1.0}. Until then, \file{baseline-v0.1} remains exactly what its name says: the comparator.

\section{Practical effort and likely payoff}

The realistic effort is 8--12 focused full-time weeks or 12--18 part-time weeks. The greatest expected gains come from data correctness, negative examples, specialist detectors, calibration, and a reliable review/export flow. A larger backbone or a fashionable new segmentation architecture is not the highest-value next step.

The likely payoff is not a spectacular single mAP number. It is a defensible demonstration that can answer:

\enlargethispage{4\baselineskip}
\begin{enumerate}[leftmargin=7mm]
  \item What sensitive content was proposed and by which provider?
  \item What did the reviewer change?
  \item How much labelled sensitive content remained visible?
  \item What ordinary content was unnecessarily hidden?
  \item Did independent detectors, OCR, barcode scanning, metadata inspection, and fresh decoding find a release problem?
  \item Can the exact result be reproduced from versioned components?
\end{enumerate}

That is a stronger final project than ``we trained Mask R-CNN and obtained 0.23 mAP.''
